\documentclass{article}

\usepackage{geometry}
\usepackage{indentfirst}
\usepackage{fancyhdr}

\pagestyle{fancy}
\geometry{a4paper, margin=1in}

\begin{document}

\title{Snapshot Objectives}
\author{Bryan Duarte}
\date{May 8th, 2026}

\maketitle

\section{Start Objective}
\fancyhf{}
\fancyhead[C]{Snapshot Objectives}
\fancyfoot[C]{\thepage}

Our start objective focuses on building the foundation of the Golden Eagle Flight Plan application. The goal is to establish the core structure of both the frontend and backend systems.

We will divide our team into frontend and backend groups to ensure efficient development. The frontend team will focus on building the user interface, while the backend team will focus on data management and API development.

Our initial features will include user authentication (login and registration), a basic dashboard layout, and database connectivity. These core systems will allow users to begin interacting with the platform.

\subsection{Front-End}
\textbf{React Native (Expo)}: Main framework used to build the mobile application. It allows for fast development and reusable UI components while maintaining performance across devices.

\textbf{Navigation System}: Provides a tabbed interface including Home, Scoreboard, Resources, AI, and Settings for easy user access.

\subsection{Back-End}

\textbf{Node.js with Express}: Backend framework used to handle API requests and manage application logic.

\textbf{MongoDB Atlas}: Cloud database used to store user data, resources, events, and progress tracking.

\section{Checkpoint 1}
At this stage, we aim to implement the core functionality of the application.

\textbf{User System}: Users can register, log in, and manage their profiles with role-based access (student or faculty).

\textbf{Scoreboard System}: A progress tracking feature that allows students to monitor achievements in Academic, Career, and Leadership categories.

\textbf{Resources and Events Feed}: Displays relevant opportunities based on user interests, major, and academic level.

\textbf{Faculty Dashboard}: Allows faculty members to post resources and events while monitoring student engagement.

\section{Checkpoint 2}
For Checkpoint 2, we plan to enhance personalization and system intelligence.

\textbf{Personalization Engine}: Matches resources and events to users using hashtags based on their profile information.

\textbf{Progress Tracking}: Tracks completed tasks and calculates levels and achievements using predefined rules.

\textbf{Security}: Implements JWT authentication to ensure secure communication and user session validation.

We also aim to refine the user interface and improve navigation for a smoother experience.

\section{Final Checkpoint}
In the final stage, we will focus on optimization and advanced features.

\textbf{Performance Optimization}: Improve loading times and overall responsiveness of the application.

\textbf{Scalability}: Ensure the system can handle more users and data efficiently.

\textbf{AI Features}: Begin integrating AI tools such as an academic advisor chatbot and resume feedback system.

\textbf{User Experience}: Enhance accessibility and usability for both students and faculty.

\section{Conclusion}
The Golden Eagle Flight Plan aims to simplify how students manage their academic and professional growth. By centralizing resources, tracking progress, and providing personalized recommendations, the system helps students stay organized and motivated. Future improvements will continue to expand AI capabilities and enhance user engagement.

\end{document}
